% 原书分章PDF第1—10页：章首导言与第26.1节。
\chapter{高阶靶向损失估计}
\label{ch:higher-order-targeted-loss}
\chapterauthors{Marco Carone、Iván Díaz 与 Mark J. van der Laan}
\chapteraffiliations{%
M. Carone（通讯作者）：美国华盛顿大学生物统计学系，F-644 Health Sciences Building, 1705 NE Pacific Street, Seattle, WA 98195, USA；电子邮箱：\href{mailto:mcarone@uw.edu}{mcarone@uw.edu}。\\
I. Díaz：美国康奈尔大学威尔康奈尔医学院卫生保健政策与研究系生物统计学及流行病学部，402 East 67th Street, New York, NY 10065, USA；电子邮箱：\href{mailto:ild2005@med.cornell.edu}{ild2005@med.cornell.edu}。\\
M. J. van der Laan：美国加利福尼亚大学伯克利分校生物统计学部及统计学系，101 Haviland Hall, \#7358, Berkeley, CA 94720, USA；电子邮箱：\href{mailto:laan@berkeley.edu}{laan@berkeley.edu}。\\
本章 DOI：\url{https://doi.org/10.1007/978-3-319-65304-4_26}。}

\begin{sourcebox}
本章旨在说明如何推广TMLE框架，使其显式利用高阶渐近表示，而不局限于一阶渐近表示。这一推广的实际意义，在于为构造有限样本表现良好、且在较弱条件下渐近有效的估计量提供指导。
\end{sourcebox}

TMLE的构造通常依赖某些一阶渐近表示。要使这些表示发挥作用，所得二阶余项必须以快于 $n^{-1/2}$ 的速率依概率趋于零。这保证一阶逼近足以指导估计量的构造，并研究其渐近极限理论。为满足这一条件，必须能够构造数据生成分布或其相关部分的一个估计量，使其收敛得足够快。例如，当数据生成分布的密度直接进入目标参数时，为保证二阶余项可以忽略，往往要求密度估计量在适当范数下以快于 $n^{-1/4}$ 的速率收敛。然而，在许多情形下，这种条件很难成立，尤其是在数据单位向量维数较高时。此时，余项本身也会影响估计量的一阶渐近行为，使其偏离渐近线性。

因此，自然会考虑依据高阶渐近展开（higher-order asymptotic expansion）构造估计量，从而计入更多解析项，并转而要求一个更高阶的余项渐近可忽略。这样，最低收敛速率条件便能显著放宽。例如，若 $r$ 阶展开存在且能够加以利用，那么，只要密度函数估计量的收敛速率快于 $n^{-1/[2(r+1)]}$，就足以使所得余项可以忽略。更为重要的是，即使这样的高阶展开并不存在，通常仍可构造近似高阶展开，从而获得明确的速率改进。不过，在后一种情形下，一般性地描述所得最低速率条件较为困难，因为这些条件可能在一定程度上依赖具体问题。

本章提出一种新的二阶TMLE（second-order TMLE），用于估计二阶路径可微的目标参数，数据来自给定半参数模型中某个未知分布的 $n$ 次独立抽样。我们将这一估计量记为2-TMLE，并以1-TMLE表示通常的一阶TMLE。通常TMLE的大样本性质要求二阶余项渐近可忽略；与此类似，这一2-TMLE的渐近正态性和有效性将依赖三阶余项的渐近可忽略性。为说明它的构造，我们针对单时点干预的g计算参数，给出一个2-TMLE。这里观测到的数据，是由基线协变量、二元处理和最终二元结局组成的随机向量的 $n$ 个独立同分布副本。为使表述简洁，本章只讨论2-TMLE，但这些技术可以直接推广至任意阶。关于这一推广的更多细节，可参见技术报告 Carone 等（2014）。

我们给出构造2-TMLE并确立其渐近有效性的一般模板。我们将说明，这一2-TMLE可以在标准TMLE框架内构造，具体做法是在标准TMLE所用的最不利参数子模型中增加一个参数。我们还将说明，这一2-TMLE的统计推断可以依据二阶Taylor展开，从而可能改善置信区间的有限样本表现。

在许多问题中，目标参数仅具有一阶路径可微性，二阶梯度并不存在。作为解决方案，我们提出基于近似二阶梯度的2-TMLE，并给出相应理论。额外的近似使所得2-TMLE多出一个偏差项，Robins 等（2009）将其称为表示误差（representation error）；2-TMLE的渐近线性与有效性定理还要求这一偏差项渐近可忽略。对这一偏差项的控制可以在标准TMLE框架内完成，而不会冒失去一阶TMLE理想性质的风险。在我们的例子中，只要适当地选择核函数与带宽速率，并且数据生成分布满足某些光滑性条件，2-TMLE的渐近有效性仍只依赖三阶余项的渐近可忽略性。我们还提出基于协同TMLE（C-TMLE）框架的数据自适应带宽选择器，其实质是在TMLE的靶向步骤中，选择使偏差减少量最大的带宽。我们通过说明如何实施2-TMLE来推断单时点干预的g计算参数，展示其用法，并报告数值结果。Díaz 等（2016）对这些内容作了详细讨论。

\section{高阶TMLE概述}

假定从 $P_0\in\mathcal M$ 得到 $n$ 个独立观测 $O_1,\ldots,O_n$，其中统计模型 $\mathcal M$ 是典型数据结构 $O$ 的所有可能概率分布所组成的集合，这些分布与关于 $P_0$ 的现有知识相容。令 $P_n$ 为 $O_1,\ldots,O_n$ 的经验分布，并定义 $\mathcal O:=\bigcup_{P\in\mathcal M}\supp(P)$，其中 $\supp(P)$ 表示 $P$ 的支撑集。令 $\Psi:\mathcal M\to\mathbb R$ 为目标参数映射，并假设我们希望根据观测数据推断真实参数值 $\psi_0:=\Psi(P_0)$。为简单起见，本章讨论单变量目标参数，但这里的所有结果也都可以直接推广至像空间包含于 $\mathbb R^d$ 的欧几里得值目标参数。

\subsection{TMLE}

这里重点考虑参数 $\Psi$ 关于 $P$ 在适当意义下足够光滑的情形。确切地说，以 $L_0^2(P)$ 表示定义在 $P$ 的支撑集上、在 $P$ 下均值为零的平方可积实值函数所成的Hilbert空间，其内积为 $\langle h_1,h_2\rangle_P:=\int h_1(o)h_2(o)\,dP(o)$。令 $T_{\mathcal M}(P)$ 为 $\mathcal M$ 在 $P\in\mathcal M$ 处的切空间，定义为经过 $P$ 的所有正则参数子模型的得分所生成的线性空间的闭包；它是 $L_0^2(P)$ 的一个子空间。若存在元素 $D^{(1)}(P)\in L_0^2(P)$，使对任意 $P_1,P_2\in\mathcal M$，线性化
\[
\begin{aligned}
\Psi(P_2)-\Psi(P_1)
 &= (P_2-P_1)D^{(1)}(P_2)+R_2(P_1,P_2)\\
 &= -P_1D^{(1)}(P_2)+R_2(P_1,P_2)
\end{aligned}
\]
成立，并且对 $\mathcal M$ 上某个距离 $\Delta$，当 $\Delta(P_1,P_2)\to0$ 时，$R_2(P_1,P_2)/\Delta(P_1,P_2)\to0$，则称 $\Psi$ 在 $\mathcal M$ 上强可微（strongly differentiable），任何这样的元素 $D^{(1)}(P)$ 都称为 $\Psi$ 在 $P$ 处的一阶梯度（Pfanzagl，1982）。这里，对任意函数 $f$ 和可能为带符号测度的 $P$，我们以 $Pf$ 表示 $\int f(o)\,dP(o)$。下文将非正式地把 $R_2(P_1,P_2)$ 称为二阶余项。

$T_{\mathcal M}(P)$ 中唯一的一阶梯度 $D^{(1)*}(P)$，既称为典范一阶梯度，也称为有效影响函数。后一术语反映了这样一个事实：$\psi_0$ 的正则估计量渐近有效，当且仅当它渐近线性，且影响函数为 $D^{(1)*}(P_0)$（Bickel 等，1997b）。形式上，$\psi_0$ 的正则估计量 $\psi_n$ 渐近有效，当且仅当
\[
\psi_n-\psi_0=(P_n-P_0)D^{(1)*}(P_0)+o_P(n^{-1/2}).
\]
TMLE为构造 $\psi_0$ 的渐近有效代入估计量 $\psi_n^*:=\Psi(P_n^*)$ 提供了一般模板。其渐近有效性，一部分源自性质
\begin{equation}
P_nD^{(1)*}(P_n^*)=\frac1n\sum_{i=1}^nD^{(1)*}(P_n^*)(O_i)=0,
\tag{26.1}
\end{equation}
以及 $\Psi$ 在 $\mathcal M$ 上的强可微性所蕴含的一阶展开
\begin{equation}
\begin{aligned}
\Psi(P)-\Psi(P_0)
 &= (P-P_0)D^{(1)*}(P)+R_2(P,P_0)\\
 &= -P_0D^{(1)*}(P)+R_2(P,P_0).
\end{aligned}
\tag{26.2}
\end{equation}
将恒等式 (26.1) 与在 $P=P_n^*$ 处计算的式 (26.2) 合并，可得
\begin{equation}
\Psi(P_n^*)-\Psi(P_0)=(P_n-P_0)D^{(1)*}(P_n^*)+R_2(P_n^*,P_0).
\tag{26.3}
\end{equation}
这是任何确立TMLE渐近线性与有效性的定理的基础。只要能够证明：
\begin{enumerate}
\item $D^{(1)*}(P_n^*)$ 以趋于一的概率属于某个 $P_0$-Donsker类；
\item $P_0[D^{(1)*}(P_n^*)-D^{(1)*}(P_0)]^2$ 依概率趋于零；
\item $R_2(P_n^*,P_0)$ 以快于 $n^{-1/2}$ 的速率依概率趋于零；
\end{enumerate}
经验过程理论的结果（例如 van der Vaart 和 Wellner，1996）便蕴含所需的渐近有效性，具体而言即
\[
\psi_n^*-\psi_0=(P_n-P_0)D^{(1)*}(P_0)+o_P(n^{-1/2}).
\]
为构造求解式 (26.1) 的估计量 $P_n^*$，TMLE依赖 $P_0$ 的一个初始估计量 $P_{n,0}$；例如，可以通过超级学习（van der Laan 和 Dudoit，2003；van der Vaart 等，2006；van der Laan 等，2007；Polley 等，2011）获得具有最优表现的初始估计量。它还依赖一个所谓的最不利参数子模型 $\{P_{n,0}(\epsilon):\epsilon\}\subset\mathcal M$，满足：
\begin{enumerate}
\item $P_{n,0}(0)=P_{n,0}$；
\item $D^{(1)*}(P_{n,0})$ 属于 $\epsilon$ 在 $\epsilon=0$ 处所有得分的线性张成空间的闭包。
\end{enumerate}

TMLE的一次更新为 $P_{n,1}:=P_{n,0}(\epsilon_n^0)$，其中
\[
\epsilon_n^0:=\argmax_{\epsilon}P_n\log p_{n,0}(\epsilon)
\]
是所构造参数子模型中的极大似然估计量（MLE），$p_{n,0}(\epsilon):=dP_{n,0}(\epsilon)/d\mu$ 是 $P_{n,0}$ 关于某个控制测度 $\mu$ 的Radon--Nikodym导数。随后，以 $P_{n,1}$ 替代 $P_{n,0}$，按照先前的方法构造最不利子模型，再重复这一过程，把 $P_{n,1}$ 映射为下一次更新 $P_{n,2}$。这一过程反复进行，直至收敛。$P_0$ 的TMLE定义为这一迭代算法的极限，记为 $P_n^*\in\mathcal M$；$\psi_0$ 的TMLE则是相应的代入估计量 $\psi_n^*:=\Psi(P_n^*)$。在每个 $\epsilon_n^j$ 为内点的步骤 $j$ 中，$\epsilon$ 的MLE $\epsilon_n^j$ 求解得分方程
% 原书行内得分方程改作无编号展示，未增加公式号。
\[
0=\left.\frac{\partial}{\partial\epsilon}P_n\log p_{n,j}(\epsilon)\right|_{\epsilon=\epsilon_n^j},
\]
其中 $p_{n,j}(\epsilon):=dP_{n,j}(\epsilon)/d\mu$ 是 $P_{n,j}$ 关于 $\mu$ 的Radon--Nikodym导数。因此，我们得到
\[
P_nD^{(1)*}(P_n^*)=
\left.\frac{\partial}{\partial\epsilon}P_n\log p_n^*(\epsilon)\right|_{\epsilon=0}=0,
\]
因为子模型 $\{P_n^*(\epsilon):\epsilon\}$ 中，$\epsilon$ 在 $\epsilon=0$ 处的得分为 $D^{(1)*}(P_n^*)$；这里，$p_n^*$ 是 $P_n^*$ 关于 $\mu$ 的Radon--Nikodym导数。

\subsection{TMLE的扩展}

上述迭代算法依赖对数似然损失，这也直接说明了靶向极大似然估计量（targeted maximum likelihood estimator）这一名称的由来。不过，只要能够找到适当的损失函数，同样的思想也可用于目标参数的任意表示 $\Psi(P)=\Psi_1(Q)$，其中 $Q:=Q(P)$ 表示计算 $\Psi(P)$ 所需的 $P$ 的相关部分，前提是 $D^{(1)*}(P):=D^{(1)*}(Q,g)$ 可以表示为 $Q$ 和干扰参数 $g:=g(P)$ 的函数。这一推广后的算法称为靶向最小损失估计（targeted minimum loss-based estimation）。它要求有 $(Q_0,g_0)$ 的初始估计量 $(Q_{n,0},g_n)$、一个损失函数 $(o,Q)\mapsto L(Q)(o)$，以及一个通常依赖于 $g_n$ 的最不利子模型 $\{Q_{n,0}(\epsilon):\epsilon\}\subset\mathcal Q$。损失函数需满足 $Q_0=\argmin_{Q\in\mathcal Q}P_0L(Q)$，其中 $\mathcal Q:=\{Q(P):P\in\mathcal M\}$；最不利子模型则需使广义得分
\[
\left.\frac{\partial}{\partial\epsilon}L(Q_{n,0}(\epsilon))\right|_{\epsilon=0}
\]
的分量张成 $D^{(1)*}(Q_{n,0},g_n)$。接着便可定义与上述过程直接对应的更新方案，以经验风险最小化替代对数似然最大化。由此得到靶向估计量 $Q_n^*$，以及 $\psi_0$ 的相应靶向最小损失估计量 $\Psi_1(Q_n^*)$，使得
\[
P_nD^{(1)*}(Q_n^*,g_n)=0,
\]
这正是渐近线性与有效性的一项基本要求。

一个简单但关键的观察是：为了赋予TMLE更多性质，可以构造 $g_0$ 的靶向估计量 $g_n^*$，使其求解指定方程，也可以很容易地调整 $Q_n^*$，使其求解其他感兴趣的方程。实现这一点，尤其可以通过在经过 $(Q_{n,0},g_n)$ 的最不利参数子模型中加入额外参数。TMLE的这一一般性此前已在多种情形下得到利用（例如 van der Laan 和 Rubin，2006；Rubin 和 van der Laan，2011；van der Laan 和 Rose，2011；Gruber 和 van der Laan，2012b；Lendle 等，2013；van der Laan，2014b），并且在本章所述高阶TMLE的构造中起着根本作用。

\subsection{二阶渐近展开}

作为TMLE基础的一阶展开 (26.2)，是参数映射 $\Psi$ 在 $\mathcal M$ 上一阶强可微的直接结果。如果所得二阶余项 $R_2(P_n^*,P_0)$ 渐近可忽略这一要求难以成立，那么转而利用二阶展开可能更为合理。这类展开所对应的余项，通常可以在较弱条件下渐近可忽略。

以 $L_0^{2*}(P^2)$ 表示定义在 $P^2$ 的支撑集上的平方可积实值函数所成的Hilbert空间。这些函数关于两个自变量对称，并对 $P$-几乎处处的 $o$ 满足
\[
\int f(o_1,o)\,dP(o_1)=\int f(o,o_2)\,dP(o_2)=0,
\]
其内积为
\[
\langle f_1,f_2\rangle_{P^2}:=P^2(f_1f_2)
=\int f_1(o_1,o_2)f_2(o_1,o_2)\,dP(o_1)\,dP(o_2).
\]
如果对任意 $P_1,P_2\in\mathcal M$，表示
\begin{equation}
\begin{aligned}
\Psi(P_2)-\Psi(P_1)
 &= (P_2-P_1)D^{(1)}(P_2)\\
 &\quad+\frac12(P_2-P_1)^2D^{(2)}(P_2)+R_3(P_1,P_2)\\
 &= -P_1D^{(1)}(P_2)-\frac12P_1^2D^{(2)}(P_2)+R_2(P_1,P_2)
\end{aligned}
\tag{26.4}
\end{equation}
对某个元素 $D^{(2)}(P_2)\in L_0^{2*}(P_2^2)$ 及某个三阶余项 $R_3(P_1,P_2)$ 成立，并且对 $\mathcal M$ 上某个距离 $\Delta$，当 $\Delta(P_1,P_2)\to0$ 时，$R_3(P_1,P_2)/\Delta^2(P_1,P_2)\to0$，则称 $\Psi$ 在 $\mathcal M$ 上二阶强可微。任何这样的元素 $D^{(2)}(P)$ 都称为 $\Psi$ 在 $P$ 处的二阶梯度（second-order gradient；Pfanzagl，1985）。

为构造2-TMLE，必须找到梯度 $D^{(2)}(P)\in L_0^{2*}(P^2)$，使对每个 $o\in\mathcal O$，$o_1\mapsto D^{(2)}(P)(o_1,o)$ 或 $o_2\mapsto D^{(2)}(P)(o,o_2)$ 属于 $T_{\mathcal M}(P)$。任何这样的元素都称为二阶部分典范梯度（second-order partial canonical gradient）。二阶得分、切空间和典范梯度的概念此前已有定义与研究（例如 Pfanzagl，1985；Robins 等，2009；van der Vaart，2014）。这些概念有助于获得二阶部分典范梯度；特别是，任何二阶典范梯度也都是二阶部分典范梯度。Carone 等（2014）对这些概念给出了自足的概述，包括强可微性与路径可微性之间的联系；我们还给出获得二阶梯度并推导二阶典范梯度的构造性方法。

对给定函数 $(o_1,o_2)\mapsto f(o_1,o_2)\in\mathbb R$，以 $f^\circ$ 表示其对称化：
\[
(o_1,o_2)\mapsto f^\circ(o_1,o_2):=\frac{f(o_1,o_2)+f(o_2,o_1)}2.
\]
注意，对任意测度 $P\in\mathcal M$，都有 $P^2f=P^2f^\circ$。正如下文将看到的，在2-TMLE的构造中，二阶部分典范梯度 $D^{(2)}(P)$ 仅通过如下形式的经验矩出现：
\[
P_n^2D^{(2)}(P),
\]
其中 $P\in\mathcal M$。因此，实践中可以使用任何元素 $D_+^{(2)}(P)\in L_0^2(P^2)$，只要对每个 $P\in\mathcal M$，都有 $P_n^2D_+^{(2)}(P)=P_n^2D^{(2)}(P)$。基于这一观察，也为简化表述，下文允许略微放宽定义：只要 $L_0^2(P^2)$ 中某元素的对称化，是 $\Psi$ 在 $P\in\mathcal M$ 处的二阶部分典范梯度或典范梯度，我们也将该元素称为二阶梯度或典范梯度。

\subsection{2-TMLE的构造}

TMLE渐近有效性的证明中，一个主要条件是二阶项 $R_2(P_n^*,P_0)$ 必须为 $o_P(n^{-1/2})$。本章的目标是构造一种TMLE，把这一二阶余项条件替换为三阶余项条件，即要求 $R_3(P_n^*,P_0)$ 为 $o_P(n^{-1/2})$。这样，便能在对TMLE初始估计量收敛速率要求更宽松、因而更现实的条件下，构造已知渐近线性且有效的估计量。

为此，需要根据 $\Psi$ 在 $P_0$ 附近的高阶渐近展开，在构造 $P_n^*$ 时加入额外的靶向拟合。如果 $\Psi$ 在 $P_0$ 处二阶强可微，即式 (26.4) 在 $P_1=P_0$ 处成立，那么，对任何满足式 (26.1) 的 $P_n^*\in\mathcal M$ 及任何二阶部分典范梯度 $D^{(2)}(P)$，都可以写出
\[
\begin{aligned}
\Psi(P_n^*)-\Psi(P_0)
 &= (P_n-P_0)D^{(1)*}(P_n^*)\\
 &\quad-\frac12P_0^2D^{(2)}(P_n^*)+R_3(P_n^*,P_0).
\end{aligned}
\]
因此，我们将在最不利参数子模型中加入额外参数，使TMLE不仅求解 $P_nD^{(1)*}(P_n^*)=0$，还求解V统计量方程
\[
0=P_n^2D^{(2)}(P_n^*)=\frac1{n^2}\sum_{i=1}^n\sum_{j=1}^nD^{(2)}(P_n^*)(O_i,O_j).
\]
这样便得到
\[
\begin{aligned}
\Psi(P_n^*)-\Psi(P_0)
 &= (P_n-P_0)D^{(1)*}(P_n^*)\\
 &\quad+\frac12(P_n^2-P_0^2)D^{(2)}(P_n^*)+R_3(P_n^*,P_0).
\end{aligned}
\]
引理26.1给出使 $(P_n^2-P_0^2)D^{(2)}(P_n^*)$ 以快于 $n^{-1/2}$ 的速率依概率趋于零的条件。接着可如前一样确立渐近有效性，只是把条件 $R_2(P_n^*,P_0)=o_P(n^{-1/2})$ 替换为 $R_3(P_n^*,P_0)=o_P(n^{-1/2})$。

我们提出的2-TMLE基于一个十分简单的观察：如果记
\[
\bar D_n^{(2)}(P_n^*)(O_i):=\frac1n\sum_{j=1}^nD^{(2)}(P_n^*)(O_i,O_j),
\]
那么 $\bar D_n^{(2)}(P_n^*)$ 本身就可以看作 $P_n^*$ 处的一个得分。这里不失一般性，假设对每个 $o\in\mathcal O$，$o_1\mapsto D^{(2)}(P)(o_1,o)$ 属于 $T_{\mathcal M}(P)$。此外，还可写成
\[
P_n^2D^{(2)}(P_n^*)=\frac1n\sum_{i=1}^n\bar D_n^{(2)}(P_n^*)(O_i)
=P_n\bar D_n^{(2)}(P_n^*).
\]
因此，我们只需把1-TMLE中在 $\epsilon_1=0$ 处得分为 $D^{(1)*}(P)$ 的最不利参数子模型 $\{P(\epsilon_1):\epsilon_1\}$，扩展为还包含参数 $\epsilon_2$ 的模型，使扩展后的最不利子模型 $\{P(\epsilon_1,\epsilon_2):\epsilon_1,\epsilon_2\}$ 在 $(\epsilon_1,\epsilon_2)=(0,0)$ 处还生成得分 $\bar D_n^{(2)}(P)$。于是，所得2-TMLE $P_n^*$ 求解方程组
\[
\begin{aligned}
0&=P_nD^{(1)*}(P_n^*),\\
0&=P_n\bar D^{(2)}(P_n^*)=P_n^2D^{(2)}(P_n^*).
\end{aligned}
\]
因此，所提出的2-TMLE就是以某种特定方式扩展了最不利子模型的普通1-TMLE。

需要注意，上述构造并不能使用任意二阶梯度。要实现TMLE构造，关键在于它的某一个坐标投影始终属于切空间 $T_{\mathcal M}(P)$。二阶典范梯度通常具有这一性质，因此通常可以据此定义适当的2-TMLE。

\subsection{可微性不足的目标参数}

正如 Robins 等（2009）所指出的，大多数感兴趣的统计参数，作为数据生成分布的函数，其光滑程度不足以对 $R_2(P_n^*,P_0)$ 作出所需展开。这些二阶余项往往包含密度估计量与真实密度之差的平方。因此，虽然 $R_2(P_n^*,P_0)$ 本身通常可以展开为 $P_n^*$ 与 $P_0$ 之间的二阶差异，但这些二阶项却不能表示为二阶梯度 $D^{(2)}(P_n^*)$ 在 $P_0$ 下的期望。这是任何旨在减少高阶偏差的方法所面临的重大障碍；尤其是，它也是构造高阶TMLE时必须面对的挑战。

在这种情形下，我们将寻找一个由平滑参数 $h$ 索引的替代对象 $D_h^{(2)}$，使对某个三阶余项 $R_3(P_n^*,P_0)$ 有
\[
R_2(P_n^*,P_0)=-\frac12\lim_{h\to0}P_0^2D_h^{(2)}(P_n^*)+R_3(P_n^*,P_0),
\]
从而得到表示
\[
R_2(P_n^*,P_0)=-\frac12P_0^2D_h^{(2)}(P_n^*)+B_n(h)+R_3(P_n^*,P_0).
\]
这里，表示误差
% 原书行内定义改作无编号展示。
\[
B_n(h):=\frac12\left[P_0^2D_h^{(2)}(P_n^*)-\lim_{h\to0}P_0^2D_h^{(2)}(P_n^*)\right]
\]
量化了近似二阶梯度所造成的偏差。因此，可写成
\[
\begin{aligned}
\Psi(P_n^*)-\Psi(P_0)
 &= (P_n-P_0)D^{(1)*}(P_n^*)-\frac12P_0^2D_h^{(2)}(P_n^*)\\
 &\quad+B_n(h)+R_3(P_n^*,P_0).
\end{aligned}
\]
为使上述2-TMLE渐近有效性的一般证明仍然成立，在对 $P_0$ 施加适当光滑性条件后，我们将要求：对适当选取的调节参数 $h=h_n$，$B_n(h)$ 以快于 $n^{-1/2}$ 的速率趋于零。当用核平滑构造近似二阶梯度时，选择趋于零足够快的带宽序列，并利用能够充分利用底层分布 $P_0$ 潜在光滑性的高阶核，通常就足以保证 $B_n(h_n)$ 渐近可忽略。

因为所用最不利子模型在 $\epsilon=0$ 处的得分同时张成 $D^{(1)*}(P)$ 和 $\bar D_{h_n}^{(2)}(P)$，我们所提2-TMLE的构造保证：对调节参数的某个取值 $h_n$，有
\[
P_nD^{(1)*}(P_n^*)=P_n^2D_{h_n}^{(2)}(P_n^*)=0.
\]
实践中，必须十分谨慎地选择调节参数 $h_n$。一方面，为确保
\[
(P_n^2-P_0^2)D_{h_n}^{(2)}(P_n^*)
\]
以快于 $n^{-1/2}$ 的速率依概率趋于零，通常要求 $h_n$ 趋于零得足够慢。另一方面，为使表示误差渐近可忽略，又要求 $h_n$ 趋于零得足够快。这构成目标参数不具二阶路径可微性时2-TMLE面临的主要理论挑战：选择 $h_n$ 时必须仔细平衡，以确保
\[
\frac12(P_n^2-P_0^2)D_{h_n}^{(2)}(P_n^*)+B_n(h_n)=o_P(n^{-1/2}).
\]
这需要在控制V统计量项与控制表示误差之间作出合理权衡。

这一理论挑战当然也直接转化为一个基本的实际挑战。实践中，必须从一大组候选2-TMLE算法中选择一个，而各个算法均由相应的 $h$ 取值索引。在具体应用中，常常可以推导 $h_n$ 的最优速率，但这些速率主要具有理论意义，对于实际如何选择 $h_n$，通常几乎不给出指导。所幸，这恰好是TMLE可以自然处理的挑战：可以根据随后参数TMLE更新步骤带来的拟合改善，评判特定调节参数 $h$ 的优劣。前面提到、后文还将讨论的C-TMLE框架，正是对此的形式化实现。
